%% ma: L10-B13-0101
%% lop: 10
%% chuong: 5
%% bai: 13
%% dang: 10.13.14
%% loai: DS
%% muc: [NB, TH, TH, VD]
%% dap_an: "ĐĐSĐ"
%% nguon: "chuyen-de-cac-so-dac-trung-cua-mau-so-lieu-khong-ghep-nhom-toan-10-chuong-trinh-moi.pdf (GV. Phan Nhật Linh) (P03), Bài 02 Dạng 2 Phần II Câu 5"
%% kien_thuc: "Lớp 10 Bài 13 (số trung bình, trung vị, tứ phân vị, mốt)"
%% trang_thai: chinh-thuc
%% ghi_chu: "Cùng số liệu Bài tập tự luận 6 (đổi tên) nhưng khác loại câu. Tự duyệt theo yêu cầu giáo viên 22/10/2026."
\begin{ex}
Bác Nam và bác Linh ghi lại số cuộc điện thoại mỗi người gọi mỗi ngày trong $10$ ngày được chọn ngẫu nhiên như sau:
\begin{center}
\begin{tabular}{|c|c|c|c|c|c|c|c|c|c|c|}\hline
Bác Nam & $2$ & $7$ & $3$ & $6$ & $1$ & $4$ & $1$ & $4$ & $5$ & $1$\\\hline
Bác Linh & $1$ & $3$ & $1$ & $2$ & $3$ & $4$ & $1$ & $2$ & $20$ & $2$\\\hline
\end{tabular}
\end{center}
Xét tính đúng sai của các khẳng định sau:
\choiceTF{\True Mốt của mẫu số liệu của bác Nam là $M_o=1$}{\True Số trung bình của số cuộc điện thoại bác Linh gọi là $3{,}9$}{Xét theo số trung bình, bác Nam gọi điện thoại thường xuyên hơn}{\True Nên dùng trung vị để so sánh số cuộc gọi của bác Nam và bác Linh}
\loigiai{a) Đúng: giá trị $1$ xuất hiện $3$ lần, nhiều nhất.

b) Đúng: $\overline x_L=\dfrac{39}{10}=3{,}9$.

c) Sai: $\overline x_N=\dfrac{34}{10}=3{,}4<3{,}9$.

d) Đúng: số liệu của bác Linh có giá trị bất thường ($20$) nên dùng trung vị ($M_e$ của bác Nam là $3{,}5$, của bác Linh là $2$).}
\end{ex}
